% =====================================================================
%  Star.org Whitepaper (English)
%  Build: npm run build:whitepaper   or   docs/whitepaper/build.sh
%  Style: preamble.tex; content mirrors docs/Star.org-Core-Philosophy.md
% =====================================================================
\documentclass[11pt,fontset=none]{ctexart}

\newcommand{\WHTITLE}{The Star.org Philosophy Whitepaper}
\newcommand{\WHTOCNAME}{Contents}
\newcommand{\WHSUBJECT}{Information Utilization as the Measure of Civilization}
\newcommand{\WHLANG}{English Edition}

% =====================================================================
%  Star.org 白皮书 · 共享排版样式
% =====================================================================
%  用法：语言源文件先 % =====================================================================
%  Star.org 白皮书 · 共享排版样式
% =====================================================================
%  用法：语言源文件先 % =====================================================================
%  Star.org 白皮书 · 共享排版样式
% =====================================================================
%  用法：语言源文件先 \input{preamble.tex}，再写正文。
%  字体使用仓库内置的 Noto Sans/Serif SC 子集（fonts/），
%  渲染不依赖运行环境是否安装中文字体（见 DECISIONS D10）。
% =====================================================================

\usepackage{fontspec}
\usepackage{xeCJK}

% ---- 字体：仓库内置的 Noto Sans/Serif SC 子集 ----
% 路径相对于编译工作目录（docs/whitepaper/）。字体随仓库内置，
% 渲染结果不再依赖运行环境是否安装中文字体（DECISIONS D10）。
% 子集覆盖 CJK 统一表意文字 U+4E00–U+9FFF 全区间（见 fonts/manifest.json）。
% 内置子集只有 Regular 一个字重，AutoFakeBold 用于合成粗体，
% 否则 \bfseries 会静默回退到 Regular、标题看不出层级。
\setCJKmainfont[Path=../../fonts/,AutoFakeBold=2.5]{NotoSerifSC-Regular.subset.otf}
\setCJKsansfont[Path=../../fonts/,AutoFakeBold=2.5]{NotoSansSC-Regular.subset.otf}
\setCJKmonofont[Path=../../fonts/]{NotoSansSC-Regular.subset.otf}
% 拉丁字母与数字沿用 TeX 自带的 Latin Modern（任何 TeX Live 环境都有），
% 不依赖系统是否安装 Helvetica 一类的西文字体。

\usepackage{geometry}
\usepackage{xcolor}
\usepackage{booktabs}
\usepackage{array}
\usepackage{tabularx}
\usepackage{longtable}
\usepackage{fancyhdr}
\usepackage{setspace}
\usepackage{parskip}
\usepackage{caption}
\usepackage{ragged2e}
\usepackage{microtype}
\usepackage[hidelinks]{hyperref}

% ---- 页面 ----
\geometry{a4paper, top=26mm, bottom=24mm, left=24mm, right=24mm}

% ---- 配色（与站点一致的深蓝 + 金）----
\definecolor{ink}{HTML}{111827}
\definecolor{inksoft}{HTML}{4B5563}
\definecolor{inkfaint}{HTML}{6B7280}
\definecolor{accent}{HTML}{1E3A8A}
\definecolor{gold}{HTML}{B08430}
\definecolor{line}{HTML}{D8DEE9}
\definecolor{surface}{HTML}{F6F7FB}

% ---- 页眉页脚 ----
\pagestyle{fancy}
\fancyhf{}
\renewcommand{\headrulewidth}{0.4pt}
\renewcommand{\footrulewidth}{0pt}
\fancyhead[L]{\footnotesize\color{inkfaint}\sffamily Star.org}
\fancyhead[R]{\footnotesize\color{inkfaint}\sffamily\WHTITLE}
\fancyfoot[C]{\footnotesize\color{inkfaint}\thepage}

% ---- 标题样式（不依赖 titlesec / sectsty）----
\makeatletter
\renewcommand{\section}{\@startsection{section}{1}{0pt}%
  {22pt plus 4pt minus 2pt}{10pt plus 2pt minus 2pt}%
  {\normalfont\sffamily\Large\bfseries\color{accent}}}
\renewcommand{\subsection}{\@startsection{subsection}{2}{0pt}%
  {16pt plus 3pt minus 2pt}{7pt plus 2pt minus 2pt}%
  {\normalfont\sffamily\large\bfseries\color{ink}}}
\renewcommand{\subsubsection}{\@startsection{subsubsection}{3}{0pt}%
  {12pt plus 2pt minus 2pt}{6pt plus 1pt minus 1pt}%
  {\normalfont\sffamily\normalsize\bfseries\color{inksoft}}}
\makeatother

% 章标题前不加首行缩进，段落间距由 parskip 控制
\setlength{\parindent}{0pt}

% ---- 引用块（核心主张）：左侧金色竖线 + 浅底 ----
\newcommand{\claim}[1]{%
  \par\vspace{6pt}%
  \noindent\begin{minipage}{\linewidth}%
    \setlength{\fboxsep}{12pt}%
    \colorbox{surface}{%
      \begin{minipage}{\dimexpr\linewidth-2\fboxsep\relax}%
        \color{accent}\bfseries\sffamily #1%
      \end{minipage}}%
  \end{minipage}\par\vspace{8pt}%
}

% ---- 一句话结语块 ----
\newcommand{\takeaway}[1]{%
  \par\vspace{8pt}\noindent\begin{minipage}{\linewidth}%
    \setlength{\fboxsep}{14pt}%
    \colorbox{accent}{%
      \begin{minipage}{\dimexpr\linewidth-2\fboxsep\relax}%
        \color{white}#1%
      \end{minipage}}%
  \end{minipage}\par\vspace{10pt}%
}

% ---- 表格统一样式 ----
\newcommand{\wtable}[1]{%
  \par\vspace{4pt}\noindent\small #1\par\vspace{6pt}%
}

% 目录条目字号
\renewcommand{\contentsname}{\sffamily\Large\bfseries\color{accent}\WHTOCNAME}

% 英文日期：ctexart 下 \today 输出中文格式（2026 年 9 月 21 日），英文版单独用这个
\newcommand{\WHdateEN}{%
  \ifcase\month\or January\or February\or March\or April\or May\or June\or
  July\or August\or September\or October\or November\or December\fi
  \space\the\day,\space\the\year}

% hyperref 书签与元数据
\hypersetup{
  colorlinks=true,
  linkcolor=accent,
  urlcolor=accent,
  pdftitle={\WHTITLE},
  pdfauthor={Star.org},
  pdfsubject={\WHSUBJECT},
}
，再写正文。
%  字体使用仓库内置的 Noto Sans/Serif SC 子集（fonts/），
%  渲染不依赖运行环境是否安装中文字体（见 DECISIONS D10）。
% =====================================================================

\usepackage{fontspec}
\usepackage{xeCJK}

% ---- 字体：仓库内置的 Noto Sans/Serif SC 子集 ----
% 路径相对于编译工作目录（docs/whitepaper/）。字体随仓库内置，
% 渲染结果不再依赖运行环境是否安装中文字体（DECISIONS D10）。
% 子集覆盖 CJK 统一表意文字 U+4E00–U+9FFF 全区间（见 fonts/manifest.json）。
% 内置子集只有 Regular 一个字重，AutoFakeBold 用于合成粗体，
% 否则 \bfseries 会静默回退到 Regular、标题看不出层级。
\setCJKmainfont[Path=../../fonts/,AutoFakeBold=2.5]{NotoSerifSC-Regular.subset.otf}
\setCJKsansfont[Path=../../fonts/,AutoFakeBold=2.5]{NotoSansSC-Regular.subset.otf}
\setCJKmonofont[Path=../../fonts/]{NotoSansSC-Regular.subset.otf}
% 拉丁字母与数字沿用 TeX 自带的 Latin Modern（任何 TeX Live 环境都有），
% 不依赖系统是否安装 Helvetica 一类的西文字体。

\usepackage{geometry}
\usepackage{xcolor}
\usepackage{booktabs}
\usepackage{array}
\usepackage{tabularx}
\usepackage{longtable}
\usepackage{fancyhdr}
\usepackage{setspace}
\usepackage{parskip}
\usepackage{caption}
\usepackage{ragged2e}
\usepackage{microtype}
\usepackage[hidelinks]{hyperref}

% ---- 页面 ----
\geometry{a4paper, top=26mm, bottom=24mm, left=24mm, right=24mm}

% ---- 配色（与站点一致的深蓝 + 金）----
\definecolor{ink}{HTML}{111827}
\definecolor{inksoft}{HTML}{4B5563}
\definecolor{inkfaint}{HTML}{6B7280}
\definecolor{accent}{HTML}{1E3A8A}
\definecolor{gold}{HTML}{B08430}
\definecolor{line}{HTML}{D8DEE9}
\definecolor{surface}{HTML}{F6F7FB}

% ---- 页眉页脚 ----
\pagestyle{fancy}
\fancyhf{}
\renewcommand{\headrulewidth}{0.4pt}
\renewcommand{\footrulewidth}{0pt}
\fancyhead[L]{\footnotesize\color{inkfaint}\sffamily Star.org}
\fancyhead[R]{\footnotesize\color{inkfaint}\sffamily\WHTITLE}
\fancyfoot[C]{\footnotesize\color{inkfaint}\thepage}

% ---- 标题样式（不依赖 titlesec / sectsty）----
\makeatletter
\renewcommand{\section}{\@startsection{section}{1}{0pt}%
  {22pt plus 4pt minus 2pt}{10pt plus 2pt minus 2pt}%
  {\normalfont\sffamily\Large\bfseries\color{accent}}}
\renewcommand{\subsection}{\@startsection{subsection}{2}{0pt}%
  {16pt plus 3pt minus 2pt}{7pt plus 2pt minus 2pt}%
  {\normalfont\sffamily\large\bfseries\color{ink}}}
\renewcommand{\subsubsection}{\@startsection{subsubsection}{3}{0pt}%
  {12pt plus 2pt minus 2pt}{6pt plus 1pt minus 1pt}%
  {\normalfont\sffamily\normalsize\bfseries\color{inksoft}}}
\makeatother

% 章标题前不加首行缩进，段落间距由 parskip 控制
\setlength{\parindent}{0pt}

% ---- 引用块（核心主张）：左侧金色竖线 + 浅底 ----
\newcommand{\claim}[1]{%
  \par\vspace{6pt}%
  \noindent\begin{minipage}{\linewidth}%
    \setlength{\fboxsep}{12pt}%
    \colorbox{surface}{%
      \begin{minipage}{\dimexpr\linewidth-2\fboxsep\relax}%
        \color{accent}\bfseries\sffamily #1%
      \end{minipage}}%
  \end{minipage}\par\vspace{8pt}%
}

% ---- 一句话结语块 ----
\newcommand{\takeaway}[1]{%
  \par\vspace{8pt}\noindent\begin{minipage}{\linewidth}%
    \setlength{\fboxsep}{14pt}%
    \colorbox{accent}{%
      \begin{minipage}{\dimexpr\linewidth-2\fboxsep\relax}%
        \color{white}#1%
      \end{minipage}}%
  \end{minipage}\par\vspace{10pt}%
}

% ---- 表格统一样式 ----
\newcommand{\wtable}[1]{%
  \par\vspace{4pt}\noindent\small #1\par\vspace{6pt}%
}

% 目录条目字号
\renewcommand{\contentsname}{\sffamily\Large\bfseries\color{accent}\WHTOCNAME}

% 英文日期：ctexart 下 \today 输出中文格式（2026 年 9 月 21 日），英文版单独用这个
\newcommand{\WHdateEN}{%
  \ifcase\month\or January\or February\or March\or April\or May\or June\or
  July\or August\or September\or October\or November\or December\fi
  \space\the\day,\space\the\year}

% hyperref 书签与元数据
\hypersetup{
  colorlinks=true,
  linkcolor=accent,
  urlcolor=accent,
  pdftitle={\WHTITLE},
  pdfauthor={Star.org},
  pdfsubject={\WHSUBJECT},
}
，再写正文。
%  字体使用仓库内置的 Noto Sans/Serif SC 子集（fonts/），
%  渲染不依赖运行环境是否安装中文字体（见 DECISIONS D10）。
% =====================================================================

\usepackage{fontspec}
\usepackage{xeCJK}

% ---- 字体：仓库内置的 Noto Sans/Serif SC 子集 ----
% 路径相对于编译工作目录（docs/whitepaper/）。字体随仓库内置，
% 渲染结果不再依赖运行环境是否安装中文字体（DECISIONS D10）。
% 子集覆盖 CJK 统一表意文字 U+4E00–U+9FFF 全区间（见 fonts/manifest.json）。
% 内置子集只有 Regular 一个字重，AutoFakeBold 用于合成粗体，
% 否则 \bfseries 会静默回退到 Regular、标题看不出层级。
\setCJKmainfont[Path=../../fonts/,AutoFakeBold=2.5]{NotoSerifSC-Regular.subset.otf}
\setCJKsansfont[Path=../../fonts/,AutoFakeBold=2.5]{NotoSansSC-Regular.subset.otf}
\setCJKmonofont[Path=../../fonts/]{NotoSansSC-Regular.subset.otf}
% 拉丁字母与数字沿用 TeX 自带的 Latin Modern（任何 TeX Live 环境都有），
% 不依赖系统是否安装 Helvetica 一类的西文字体。

\usepackage{geometry}
\usepackage{xcolor}
\usepackage{booktabs}
\usepackage{array}
\usepackage{tabularx}
\usepackage{longtable}
\usepackage{fancyhdr}
\usepackage{setspace}
\usepackage{parskip}
\usepackage{caption}
\usepackage{ragged2e}
\usepackage{microtype}
\usepackage[hidelinks]{hyperref}

% ---- 页面 ----
\geometry{a4paper, top=26mm, bottom=24mm, left=24mm, right=24mm}

% ---- 配色（与站点一致的深蓝 + 金）----
\definecolor{ink}{HTML}{111827}
\definecolor{inksoft}{HTML}{4B5563}
\definecolor{inkfaint}{HTML}{6B7280}
\definecolor{accent}{HTML}{1E3A8A}
\definecolor{gold}{HTML}{B08430}
\definecolor{line}{HTML}{D8DEE9}
\definecolor{surface}{HTML}{F6F7FB}

% ---- 页眉页脚 ----
\pagestyle{fancy}
\fancyhf{}
\renewcommand{\headrulewidth}{0.4pt}
\renewcommand{\footrulewidth}{0pt}
\fancyhead[L]{\footnotesize\color{inkfaint}\sffamily Star.org}
\fancyhead[R]{\footnotesize\color{inkfaint}\sffamily\WHTITLE}
\fancyfoot[C]{\footnotesize\color{inkfaint}\thepage}

% ---- 标题样式（不依赖 titlesec / sectsty）----
\makeatletter
\renewcommand{\section}{\@startsection{section}{1}{0pt}%
  {22pt plus 4pt minus 2pt}{10pt plus 2pt minus 2pt}%
  {\normalfont\sffamily\Large\bfseries\color{accent}}}
\renewcommand{\subsection}{\@startsection{subsection}{2}{0pt}%
  {16pt plus 3pt minus 2pt}{7pt plus 2pt minus 2pt}%
  {\normalfont\sffamily\large\bfseries\color{ink}}}
\renewcommand{\subsubsection}{\@startsection{subsubsection}{3}{0pt}%
  {12pt plus 2pt minus 2pt}{6pt plus 1pt minus 1pt}%
  {\normalfont\sffamily\normalsize\bfseries\color{inksoft}}}
\makeatother

% 章标题前不加首行缩进，段落间距由 parskip 控制
\setlength{\parindent}{0pt}

% ---- 引用块（核心主张）：左侧金色竖线 + 浅底 ----
\newcommand{\claim}[1]{%
  \par\vspace{6pt}%
  \noindent\begin{minipage}{\linewidth}%
    \setlength{\fboxsep}{12pt}%
    \colorbox{surface}{%
      \begin{minipage}{\dimexpr\linewidth-2\fboxsep\relax}%
        \color{accent}\bfseries\sffamily #1%
      \end{minipage}}%
  \end{minipage}\par\vspace{8pt}%
}

% ---- 一句话结语块 ----
\newcommand{\takeaway}[1]{%
  \par\vspace{8pt}\noindent\begin{minipage}{\linewidth}%
    \setlength{\fboxsep}{14pt}%
    \colorbox{accent}{%
      \begin{minipage}{\dimexpr\linewidth-2\fboxsep\relax}%
        \color{white}#1%
      \end{minipage}}%
  \end{minipage}\par\vspace{10pt}%
}

% ---- 表格统一样式 ----
\newcommand{\wtable}[1]{%
  \par\vspace{4pt}\noindent\small #1\par\vspace{6pt}%
}

% 目录条目字号
\renewcommand{\contentsname}{\sffamily\Large\bfseries\color{accent}\WHTOCNAME}

% 英文日期：ctexart 下 \today 输出中文格式（2026 年 9 月 21 日），英文版单独用这个
\newcommand{\WHdateEN}{%
  \ifcase\month\or January\or February\or March\or April\or May\or June\or
  July\or August\or September\or October\or November\or December\fi
  \space\the\day,\space\the\year}

% hyperref 书签与元数据
\hypersetup{
  colorlinks=true,
  linkcolor=accent,
  urlcolor=accent,
  pdftitle={\WHTITLE},
  pdfauthor={Star.org},
  pdfsubject={\WHSUBJECT},
}


\begin{document}

% ============================ Cover ============================
\begin{titlepage}
\centering
\vspace*{28mm}

{\sffamily\footnotesize\color{gold}STAR.ORG}\\[6mm]

{\sffamily\Huge\bfseries\color{accent} The Core Philosophy}\\[5mm]

{\sffamily\large\color{inksoft} Information Utilization as the Measure of Civilization}\\[16mm]

\textcolor{line}{\rule{\linewidth}{0.6pt}}\\[8mm]

{\normalsize\color{inksoft}
Humanity already measures the height of a civilization by its energy.\\
We want to establish how much of the universe we \emph{understand} as a measure in its own right,\\
and we believe the marks on this scale should be made by every ordinary person, together.}\\[8mm]

\textcolor{line}{\rule{\linewidth}{0.6pt}}

\vfill

{\small\color{inkfaint}
Version: v1.1 \quad·\quad \WHLANG \quad·\quad \WHdateEN\\[2mm]
Companion document: \texttt{docs/Star.org-Core-Philosophy.md}\\[1mm]
Chinese Edition: \texttt{star-org-whitepaper-zh.pdf}}

\vspace*{16mm}
\end{titlepage}

% ============================ Contents ============================
\tableofcontents
\newpage

% ============================ Abstract ============================
\section*{Abstract}
\addcontentsline{toc}{section}{Abstract}

This paper sets out a framework for thinking about civilization and the development of space, in five core ideas:

\begin{enumerate}
  \item \textbf{A civilization's level should be defined by its ``information utilization,''} not only by its ``energy utilization.'' Measured this way, Earth is currently at \textbf{Information Civilization, Level One} --- on the criterion that \textbf{observation and action have not yet closed the loop}, not on the idea that ``we know little'' (see 1.4).
  \item The universe contains three fundamentally different resources: \textbf{information ore, energy ore, and matter ore}. Energy ore and matter ore demand enormous industrial and spacefaring capacity; information ore demands only observation and attention. Star.org focuses on developing \textbf{information ore}.
  \item To mobilize all of humanity to develop space, each person must first forge a real bond with the sky; and the key to that bond is \textbf{the right to name}.
  \item Whoever observes or attends to a star's information is developing that star's information ore, and therefore holds a certain \textbf{information right} over it.
  \item Our bodies cannot reach the distant stars, but our \textbf{information and imagination} can. Science fiction, film, and games are likewise a mode of mining information ore.
\end{enumerate}

The five ideas are not parallel slogans but a single chain of reasoning that joins end to end: from the yardstick, to the resource, to the participants, to the rights, and finally to another way of reaching the cosmos.

\vspace{4mm}
\takeaway{The height of a civilization is decided not only by how much energy it can harness, but also by how much of the universe it can see clearly.}

{\small\color{inkfaint}New in v1.1: section 1.4, ``Why `Level One': the observation--action loop criterion'', which turns the level from a matter of rhetoric into a recomputable criterion; and section 1.5, ``The level is not a scalar: at least four axes'', a self-assessment along coverage, depth, participation and closure.\par}

\newpage

% ============================ Body ============================
\section{Civilization Level: from Energy Utilization to Information Utilization}

\subsection{The Reference Frame: Existing Scales Are Defined by Energy}

The most commonly cited yardstick for civilization is \textbf{energy utilization} — how much energy a civilization can harness determines its place on the cosmic ladder. This treats a civilization as ``a machine of ever-greater power'': from burning wood and coal, to nuclear fission, to the hypothetical Dyson sphere that wraps a star. The scale is \textbf{a magnitude of energy}.

This yardstick carries a hidden premise: \textbf{it measures what a civilization can \emph{do}, not what it \emph{knows}.}

\subsection{Star.org's Yardstick: Information Utilization}

Star.org proposes another scale that is equally fundamental yet often overlooked — \textbf{information utilization}:

\claim{A civilization's information utilization = the degree to which it grasps true information about the universe, and its capacity to turn that information into understanding and action.}

In plain terms: \textbf{how much of the universe we have cognitively covered, and how deeply we have dug into it}. It measures a civilization's cognitive frontier, not its power ceiling.

The two scales are not mutually exclusive but complementary: energy determines whether a civilization can move; information determines whether it can see. Star.org chooses to center on the information side — because it is closer to ``people,'' and closer to the part that every ordinary person can take part in.

\subsection{Earth Today: Information Civilization, Level One}

Measured against this scale, Star.org gives Earth's present stage a designation:

\claim{Earth is currently at ``Information Civilization, Level One.''}

``Level One'' means something quite plain: we have only just begun to \textbf{systematically and planet-wide} record and share information about the universe. The celestial bodies we can see, measure, and describe remain an extremely small patch relative to the whole cosmos; and our grasp of that information mostly stays at the level of ``knowing its position, brightness, and distance.'' \textbf{This is the starting point of information utilization — not its endpoint.}

Calling Earth ``Information Civilization, Level One'' is not a put-down but an attempt to supply \textbf{a coordinate system one can climb} — so that we know exactly which rung we stand on, and what the next rung depends on.

\subsection{Why ``Level One'': the observation--action loop criterion}

``Level One'' should not be a humble phrase; it should be a criterion \textbf{anyone can recompute}: whether a civilization's \textbf{reach of observation and reach of action are symmetric}.

Put the two kinds of ``reach'' side by side:

\wtable{
\renewcommand{\arraystretch}{1.35}
\begin{tabularx}{\linewidth}{@{}p{34mm} X@{}}
\toprule
\textbf{Direction} & \textbf{Farthest reach} \\
\midrule
\textbf{Information} (the light we receive) & The oldest light comes from about 13.8 billion years ago; the observable universe has a comoving radius of about \textbf{46 billion light-years} \\
\textbf{Matter} (the things we send out) & Voyager 1 is about \textbf{172 AU} (roughly 0.0027 light-years), about \textbf{0.06\%} of the way to the nearest star (Proxima Centauri, about 4.24 light-years) \\
\textbf{Their ratio} & about \textbf{13 orders of magnitude} \\
\bottomrule
\end{tabularx}}

Our eyes already reach the edge of the universe; our hands are still in our own backyard. \textbf{Seeing without being able to act is the defining feature of Level One.}

That also lets us write down the threshold for Level Two:

\claim{Level Two = observation and action have closed the loop (information reach $\approx$ action reach), or the civilization has formed a self-correcting, systematic cognitive loop over the universe's information.}

\subsection{The level is not a scalar: at least four axes}

Compressing ``information utilization'' into a single number makes it easy to abuse — whichever axis you want to stress, that is the conclusion you get. So it should be split into at least four axes that can be discussed separately:

\wtable{
\renewcommand{\arraystretch}{1.35}
\begin{tabularx}{\linewidth}{@{}p{21mm} X p{44mm}@{}}
\toprule
\textbf{Axis} & \textbf{Meaning} & \textbf{Earth today (approx.)} \\
\midrule
Coverage & how much of the universe we have sampled & Very low: Gaia has measured about 1.8 billion stars, while the Milky Way holds hundreds of billions; about 6{,}400 exoplanets are confirmed against hundreds of billions of stars (about $10^{-8}$) \\
Depth & how well we understand what we sampled & Moderate to high: QED's electron magnetic moment matches experiment to more than ten significant digits, and the universe's age and main composition are measured to around one percent — yet 95\% of the energy density still has only names (dark matter, dark energy) \\
Participation & how many people actually use this information & Moderate: about 6 billion people can go online and query astronomical databases, but very few ever look \\
Closure & whether observation turns into action & Not formed: see 1.4 \\
\bottomrule
\end{tabularx}}

Writing it as a vector rather than a number guards against self-deception: \textbf{when we say ``Level One'', we mean ``the loop is not closed'', not ``we know little.''} By depth alone we are not low; by coverage alone we are near zero — both can be true at once, and any single number must hide one of them.

{\small\color{inkfaint}
On the figures: Voyager 1's distance is the 2026 measured value; the observable universe's comoving radius of about 46.1--46.6 billion light-years follows from Planck cosmological parameters; the exoplanet count is NASA Exoplanet Archive data for September 2026; the internet-user figure is the ITU 2025 estimate. Such numbers are updated as observations improve — readers are welcome to recompute them.\par}

\section{Three Kinds of Ore in the Universe: Information, Energy, Matter}

If we treat the universe as a treasure to be developed, Star.org holds that it consists of three fundamentally different kinds of resources:

\wtable{
\renewcommand{\arraystretch}{1.35}
\begin{tabularx}{\linewidth}{@{}p{21mm} X p{32mm} p{13mm}@{}}
\toprule
\textbf{Resource} & \textbf{What it is} & \textbf{Who can develop it} & \textbf{Here} \\
\midrule
Energy ore & The energy of the universe that can be harnessed (radiation, gravity, fusion, at the scale of total quantities) & Civilizations with engineering and industrial capacity & No \\
Matter ore & Physical substance in the universe (planets, ore, water, gas — collectable or manufacturable) & Civilizations able to reach and extract & No \\
Information ore & \textbf{Information about celestial bodies and phenomena} (position, magnitude, distance, spectrum, motion, history, even the stories behind them) & Requires only \textbf{observation and attention} — open to all & \textbf{Yes} \\
\bottomrule
\end{tabularx}}

The difference among the three is not ``how precious'' they are, but \textbf{the barrier to entry}:

\begin{itemize}
  \item Developing energy ore and matter ore both requires a civilization to first build \textbf{enormous industrial and spacefaring capacity} before it can get its hands on the stuff. For an individual, that is out of reach.
  \item \textbf{The barrier to information ore is far lower}: what it demands is not a starship but \textbf{attention}. An ordinary person need only be willing to look up, to observe, to care about a particular star — and he is already participating in the development of the universe's information ore.
\end{itemize}

\claim{Star.org focuses on developing information ore — because among the three, it is the only resource that lets all of humanity get started immediately, without first colonizing the cosmos.}

\section{How to Mobilize All of Humanity: a Bond with the Stars, and the Right to Name}

\subsection{The Biggest Bottleneck Is That Space Feels Unrelated to Everyone}

The grand narrative of humanity developing the cosmos tends to make ordinary people feel that ``this is a matter for astronauts, scientists, and billionaires.'' When something has no concrete connection to an individual, it generates no real drive to participate. Star.org believes:

\claim{To get all of humanity to develop space, the first thing is to give ``the universe'' and ``each person'' a real, tangible bond — rather than letting the universe remain a distant background that has nothing to do with them.}

\subsection{The Key to That Bond: the Right to Name}

Among all the ways of forging a bond, Star.org holds that the most fundamental and most moving is —

\claim{the right to name.}

To name a star is humanity's oldest and most instinctive act of cognition. From ancient mariners navigating by the constellations, to astronomers numbering celestial bodies, \textbf{``what you call it'' has always been the earliest bond forged between people and the stars.} Once a name is given, that star is no longer a point of light on a backdrop; it becomes ``my star'' — a concrete, identifiable, callable sense of belonging.

Star.org treats this ``right to name'' as the \textbf{key} connecting people to the universe: when a person names and claims a real, existing star, an \textbf{irreplaceable, personal bond} comes into being between them and that patch of sky. Developing space is no longer only a grand slogan; it becomes something concrete — ``the star I claimed is being recorded, being watched, being seen by others.''

\subsection{An Important Clarification: This Is Not Official Astronomical Naming}

Star.org fully respects the authority of the astronomical community and the legitimate systems of astronomical naming. ``Naming/claiming'' here is a \textbf{personal bond of emotion and cognition} — an \textbf{observation and memorial marker} that a namer establishes for a real celestial body. It \textbf{does not constitute, and does not claim to be,} a naming right granted by the International Astronomical Union (IAU) or any other official body.

This clarification matches the site: every page must carry the ``not official IAU naming'' statement.

\subsection{From This, Everyone's Effort Gathers}

When everyone can forge a bond with the universe through naming, developing space ceases to be the enterprise of a few and becomes instead:

\begin{itemize}
  \item each person \textbf{paying attention} to a patch of sky and recording what they observe;
  \item each person \textbf{sharing} that bond, so that more people are willing to look up;
  \item the \textbf{attention} of countless individuals converging — which is itself the broadest and most sustainable development of the information ore.
\end{itemize}

\claim{Mobilizing all of humanity to develop space does not rely on commands. It relies on giving every person a star of their own.}

\section{Information Rights: Whoever Attends, Develops — and Therefore Holds}

\subsection{The Core Claim}

Building on the definition of developing information ore, Star.org puts forward a claim about \textbf{information rights}:

\claim{Whoever observes or pays attention to a star's information is developing that star's information ore; and this act of development gives that person a certain kind of right over that information.}

Unpacked, the logic runs in three steps:

\begin{enumerate}
  \item \textbf{Attention = development}: for a star's information ore, development does not mean digging it open but \textbf{observing, recording, understanding, and telling} it. Every act of attention is an act of extraction from that star's information ore.
  \item \textbf{Development = contribution}: this information is not anyone's private possession, but the \textbf{attention, observation, stories, and memory} a person invests in it are real and indelible contributions.
  \item \textbf{Investment = right}: since ``whoever develops it holds it'' is the general principle of resource development, then \textbf{whoever pays attention to a star holds a certain right over its information} — a right that can be marked, recorded, preserved long-term, and continued for as long as that person exists.
\end{enumerate}

\subsection{What This Right Is, and Is Not}

\wtable{
\renewcommand{\arraystretch}{1.35}
\begin{tabularx}{\linewidth}{@{}X X@{}}
\toprule
\textbf{It is} & \textbf{It is not} \\
\midrule
A bond of attention that is \textbf{identifiable, recordable, and preservable} & \textbf{Physical belonging} of the star itself \\
A \textbf{record of development} over this star's \textbf{information ore} & Exclusive possession of all astronomical data about a star (the real data belongs to the scientific commons and to all humanity) \\
A \textbf{memorial and emotional belonging} that belongs to the namer & A transferable, appreciable \textbf{financial instrument} or an exchange-listed object \\
\bottomrule
\end{tabularx}}

The site sells no ownership, makes nothing into a financial instrument, and points to no investment or appreciation — this is the compliance red line, enforced automatically by \texttt{scripts/check-compliance.mjs}.

\subsection{Publicly Verifiable: Making the Right Visible and Checkable}

Star.org holds that a right is meaningful only if \textbf{others can independently verify it.} Therefore every claimed star corresponds to:

\begin{itemize}
  \item a \textbf{public claim record} (containing the star's real astronomical data);
  \item a \textbf{permanent link} that anyone can visit and verify at any time;
  \item a \textbf{certificate} that can be preserved long-term.
\end{itemize}

Anyone who wants to verify ``whether this star has been claimed, and by whom'' can check for themselves in the public claim registry. This \textbf{openness, transparency, and traceability} is precisely the foundation on which information rights can stand — it turns a private act of attention into a \textbf{public record} that others can see and acknowledge.

\section{We Cannot Reach the Stars in Body, but We Can in Information and Imagination}

\subsection{A Fundamental Constraint: We Cannot Reach the Stars}

Physically, the human body is confined to Earth. We cannot stand beside another planet or another nebula and mine it in person — at least for any foreseeable future, \textbf{energy ore and matter ore are beyond the reach of the vast majority of people.} This is the coldest fact about developing space.

\subsection{But ``We'' Are More Than a Body}

Star.org holds that what defines a person, or a civilization, is not only where its body resides, but \textbf{how far its information has reached and how far its imagination extends.}

\claim{We cannot reach the distant stars in body, but we can in our information and our imagination.}

These two things — \textbf{information and imagination} — reach deep into the cosmos without a starship. They are the true extension of humanity's territory: through imagination and storytelling, humanity has long since moved into the stars in spirit.

\subsection{Therefore, the Production of Imagination Is Itself a Form of Development}

Since the development of information ore depends on attention and cognition rather than physical arrival, then every \textbf{work that directs human attention and imagination toward the cosmos} is mining information ore:

\begin{itemize}
  \item \textbf{Science fiction}: \emph{The Three-Body Problem} directs human imagination toward the dark forest of the cosmos, interstellar civilizations, and questions of scale;
  \item \textbf{Science-fiction film and television}: \emph{Star Trek} made ``to explore the unknown'' a conviction for generations;
  \item \textbf{Games}: letting players sail, colonize, and encounter civilizations in virtual starfields;
  \item \textbf{Popular astronomy, documentaries, and art}: translating real celestial information into stories ordinary people can feel.
\end{itemize}

\textbf{None of these are escapes from reality; they are a special mode of mining the universe's information ore} — they extract not a gram of matter and not a wisp of energy, yet they mine the most precious thing of all: \textbf{humanity's understanding and imagination of who we are and where in the universe we stand.}

\subsection{The Mining of Imagination and the Mining of Real Data Complete Each Other}

Imagination is not fabricated from nothing; it needs \textbf{real celestial information} as its soil. And real celestial information needs imagination to give it meaning and draw more people to attend to it.

\begin{itemize}
  \item Science fiction inspires ordinary people to \textbf{attend to real stars};
  \item Those who attend to real stars supply \textbf{new imaginative material} for science fiction and for science.
\end{itemize}

Thus \textbf{the sky of \emph{The Three-Body Problem}, the interstellar federation of \emph{Star Trek}, and the star charts of a game, together with the coordinates, distances, and spectra of real celestial bodies, form one superimposed information ore} — all of it mined by human attention and imagination, and all of it expanding humanity's information utilization.

\subsection{What This Means for Star.org}

This explains why Star.org both handles the claiming and recording of real celestial bodies and can accommodate sky narratives at the level of fantasy:

\begin{itemize}
  \item A person claiming a \textbf{real} star is mining that star's information ore;
  \item A person stirred toward the sky by science fiction, a game, or a story is likewise \textbf{mining humanity's collective information ore.}
\end{itemize}

\claim{The two are not opposed; they are different veins of the same information ore.}

Star.org's standpoint is that every vein should point in the same direction: \textbf{to turn more people's attention and imagination toward that real sky overhead — traceable, verifiable, and above them all.}

\section{How the Five Ideas Form One Line}

\begin{center}
\sffamily\small
\begin{tabularx}{0.96\linewidth}{@{}l X@{}}
\toprule
\textbf{Yardstick} & Energy utilization → \textbf{information utilization} (Star.org's measure) \\
\textbf{Resource} & The universe has three ores: information / energy / matter \\
\textbf{Choice} & Star.org focuses on information ore — the only one open to all \\
\textbf{Participation} & Developing information ore requires all of humanity \\
\textbf{Key} & The key to bonding each person with a star = the right to name \\
\textbf{Rights} & Whoever attends = whoever develops = whoever holds its information right \\
\textbf{Extension} & Not in body, but in information and imagination: \emph{The Three-Body Problem} / \emph{Star Trek} / games \\
\bottomrule
\end{tabularx}
\end{center}

\vspace{4mm}
Every act of attention plus every act of imagination converges into humanity's shared advance in understanding the universe.

\section{Conclusion}

\takeaway{%
The height of a civilization is decided not only by how much energy it can harness, but also by how much of the universe it can see clearly.\\[2mm]
There are three ores in the universe; Star.org mines the information ore — and the mining of information ore belongs to everyone willing to look up, to name a star, and to attend to it for a long time.\\[2mm]
Whoever attends, develops; whoever develops, holds — this information right, we make public, verifiable, and belonging to every person.\\[2mm]
We cannot reach the distant stars in body, but we can in our information and imagination — The Three-Body Problem, Star Trek, and every starfield of every game are mining the information ore for humanity.\\[3mm]
Claim a star, then keep observing it — together we raise humanity's information civilization level in the universe.}

\vspace{6mm}
\noindent\textcolor{line}{\rule{\linewidth}{0.6pt}}\\[3mm]
{\small\color{inkfaint}
Star.org provides a commemorative, publicly verifiable star claiming service. It does not constitute official naming by the International Astronomical Union (IAU), and does not represent any celestial naming right. The astronomical data cited in this whitepaper comes from public star catalogues.}

\end{document}
